\section{The functional \texorpdfstring{$Y$}{Y}: Fourier truncation and estimates preserving the density}
\label{sec:lifts}

\paragraph{What this section proves and how to read it.}
This section defines $Y=\Yf(\Lsc,A)$, the $72$-fold average of a
truncated Fourier series $F_A$ of $\one_A$ (the \emph{lift}: $F_A$
retains characters of conductor at most $N^{1/20}$ so that fixing
labels can be analysed prime by prime), and proves three properties
used throughout: averaging interval-and-residue inputs returns $F_A$
exactly; $L^{2r}$ norms of lifts are linear in the density $\alpha$
(so remainders in \cref{sec:scale-comparison} are $O(\Lsc^{-\sigma}
\alpha)$ with $\sigma=1/825$ (\cref{def:record:analytic-stage})); and Fourier terms of large conductor are damped at large
primes. It also proves that tensor products preserve exact Fourier
support (the resulting orthogonal components of a function,
indexed by exact Fourier support, are called \emph{sectors}; they
are not prime pieces; a \emph{spectator} is a factor left
unchanged in such a tensor product). Inputs: \cref{sec:positivity} and the reference hypotheses
(R1)--(R5) (\cref{subsec:lifts:reference}), verified in Part~V.
Outputs: the definition \cref{def:record:functional} and the moment
and decay estimates cited in Parts~III and~V. A first reading should
follow the definition and the statements of the moment and decay
results; the tensor-product linear algebra can be used as reference.
(R4) uses \emph{mixing} (damping of nonconstant functions
at large primes, \cref{def:mixing} below) and (R5) uses the forcing terminology
of \cref{subsec:intro:glossary}.

\subsection{Definition of the lift \texorpdfstring{$F_A$}{F\_A} and the functional \texorpdfstring{$Y(\Lsc,A)$}{Y(L,A)}}\label{subsec:lifts:functional}

The slot array is $V=(\Z/6\Z)^2\times\{0,1\}$ of \cref{def:slot-array},
$d=|V|=72$, with the family $\mathfrak F$ of its cuts: the chamber cuts
$\vartheta_{c,a}$ of the two chamber coordinates and the reserved cut $\iota$.
The family contains a sorting sequence (\cref{def:sorting},
\cref{lem:slot-array}(iii)). The multilinear averages
$\Lambda_\kappa((g_v)_v)$ and $\Lambda_\kappa(g)$ of a finite signed measure
$\kappa$ on $X^V$ are those of \cref{sec:positivity}.

\begin{definition}[The functional]\label{def:record:functional}
Let $\beta=1/20$, let $\Kint\ge2$ be a fixed integer, and let $Q_3=19951586$,
$\Pfour=2^{512}$ and $\Pfive=2^{65536}$. Put
\[
 \mathcal P_1=\{p<\Pfour:\ p\equiv1\ (\mathrm{mod}\ 4)\},\qquad
 \mathcal P_3=\{Q_3\le p<\Pfour:\ p\equiv3\ (\mathrm{mod}\ 4)\},
\]
fix integers $m_p\ge1$ for $p\in\mathcal P_1$, and put $e_p=2m_p$ for
$p\in\mathcal P_1$ and $e_p=1$ for every other prime. For a finite set
$\mathcal S$ of primes let $X_{\mathcal S}=\prod_{p\in\mathcal S}\Z/p^{e_p}\Z$,
with its uniform probability measure, and call
$D_{\mathcal S}=|X_{\mathcal S}|=\prod_{p\in\mathcal S}p^{e_p}$ its \emph{full
order}.
\begin{enumerate}[label=(\alph*)]
\item \emph{Scales.} For an integer $\Lsc\ge1$ put $N=\Kint\Lsc$, $Q=N^\beta$,
 $\mathcal T_N=\{p\ \text{prime}:\ \Pfour\le p\le N^\beta\}$,
 $\mathcal P_N=\mathcal P_1\cup\mathcal P_3\cup\mathcal T_N$ and
 $G_N=X_{\mathcal P_N}$.
\item \emph{Characters.} A character of $X_{\mathcal S}$ is
 $z\mapsto e(z\cdot\xi)$ with $\xi=\sum_{p\in\mathcal S}a_p/p^{e_p}\in\Q/\Z$,
 $a_p\in\Z/p^{e_p}\Z$, and $z\cdot\xi=\sum_pa_pz_p/p^{e_p}$; we identify it
 with $\xi$. Its \emph{exact support} is $\{p:a_p\ne0\}$, and its
 \emph{conductor} $q(\xi)$ is the reduced denominator of $\xi$, that is,
 $q(\xi)=\prod_pp^{e_p-\min(v_p(a_p),e_p)}$, the product over
 $p\in\mathcal S$, with $a_p$ represented by an integer between $0$ and
 $p^{e_p}-1$ and $v_p(0)$ taken as $+\infty$. For $n\in\Z$ the number
 $e(n\xi)$ is well defined. The Fourier coefficient of
 $G\colon X_{\mathcal S}\to\C$ is $\widehat G(\xi)=\E_zG(z)e(-z\cdot\xi)$.
\item \emph{Lift.} For a nonempty interval $I$ of integers and
 $f\colon I\to\C$ put
 \[
  \widehat f_I(\xi)=\frac1{|I|}\sum_{n\in I}f(n)e(-n\xi)\quad(\xi\in\Q/\Z),
  \qquad
  \mathcal L_{I,N}f(z)=\sum_{\xi\in\Xi_N}\widehat f_I(\xi)\,e(z\cdot\xi)\quad(z\in G_N),
 \]
 where $\Xi_N$, the set of \emph{retained characters}, consists of the
 characters $\xi$ of $G_N$ with $q(\xi)\le Q$. It is defined by the
 \emph{total} conductor and depends neither on $I$ nor on $f$. It contains
 every character supported on the prefix coordinates ($p\in\mathcal P_1$) of conductor $p^j\le Q$, $1\le j\le e_p$, and it is
 closed under negation, so the lift of a real function is real.
\item \emph{Atoms and reference laws.} The set $\mathcal A_N$ of \emph{atoms}
 at scale $N$ consists of the clique atoms $\{p\}$, $p\in\mathcal P_1$; the
 blocks of a fixed partition of $\mathcal P_3$ into finite sets (the middle
 blocks of \cref{def:record:exact-blocks} and the upper blocks of
 \cref{def:record:upper-bands}); and the tail atoms $\{p\}$,
 $p\in\mathcal T_N$. For an atom $\Blk$ we write $X_\Blk$ for $X_{\mathcal S}$
 with $\mathcal S$ the set of primes of $\Blk$, and $\mu_\Blk$ for its
 reference law, a probability measure on $X_\Blk^V$ (constructed in Part~V;
 Parts~II and~III use only \cref{hyp:record:reference}). For a set
 $\mathcal S$ of atoms put $\mu_{\mathcal S}=\bigotimes_{\Blk\in\mathcal S}\mu_\Blk$
 and $Z_{\mathcal S}(g)=\Lambda_{\mu_{\mathcal S}}(g)$ for
 $g\colon X_{\mathcal S}\to\R$; and $\mu_N=\mu_{\mathcal A_N}$,
 $Z_N=Z_{\mathcal A_N}$.
\item \emph{Functional.} For $A\subseteq[N]$ let $F_A=\mathcal L_{[N],N}\one_A$
 and $\Yf(\Lsc,A)=Z_N(F_A)$, and let $\Ym(\Lsc)$ be the maximum of
 $\Yf(\Lsc,A)$ over the square-difference-free sets $A\subseteq[N]$.
\item \emph{Inputs.} For $b\in\Z/\Kint\Z$, identified with
 $\{1,\dots,\Kint\}$ through the least positive residues, let
 $I_b=((b-1)\Lsc,b\Lsc]\cap\Z$, and put $\mathcal I_\Lsc=\{[N],I_1,\dots,I_\Kint\}$.
 The intervals $I_b$ have length $\Lsc$ and partition $[N]$; every element of
 $I_b$ is smaller than every element of $I_{b+1}$ for $1\le b<\Kint$, and the
 shift $b\mapsto b+1$ wraps $I_\Kint$ to $I_1$. The \emph{small modulus} is
 $\Mzero=2^{2m_{\rm ext}+1}\prod_{3\le p<Q_3,\ p\equiv3\ (4)}p^{2m_{\rm ext}}$, with fixed $m_{\rm ext}=64$. For $I\in\mathcal I_\Lsc$, a
 positive divisor $D_0$ of $\Mzero$ and $s\in\Z/D_0\Z$,
 \[
  g^{I,D_0,s}=\mathcal L_{I,N}\big((D_0\one_{n\equiv s\,(D_0)}\one_A)|_I\big).
 \]
\end{enumerate}
\end{definition}

Throughout, $Y$ and $\Yf$ denote this functional. The small modulus is coprime
to $D_{\mathcal P_N}$: its primes are $2$ and the primes $p\equiv3\pmod4$ below
$Q_3$, none of which lies in $\mathcal P_N$. So $\Mzero$ is not a coordinate of
$G_N$; it enters only through the inputs and the external index of the
queries (\cref{def:hyb:coupled}). The atoms of $\mathcal P_1\cup\mathcal P_3$
are present at every scale, also when their primes exceed $N^\beta$; a clique
prime $p>Q$ then carries no nonconstant retained character. Two fixed integers
control them,
\[
 Q_{\rm cl}=\prod_{p\in\mathcal P_1}p^{e_p},\qquad
 Q_{\rm pre}=Q_{\rm cl}\prod_{\substack{Q_3\le p<\Pfive\\ p\notin\mathcal P_1}}p ,
\]
the full order of the clique prefix and a common multiple of the full orders of
all atoms below $\Pfive$. The data $\Kint$ and $(m_p)_{p\in\mathcal P_1}$ are
fixed, and constants may depend on them and on $\Mzero$, $Q_{\rm cl}$,
$Q_{\rm pre}$ (\cref{rem:fixed-parameters}). The baseline uses
$\Kint=\PushK$, and the main theorem uses $\Kint=\FourRoutesK$;
both use $m_p=\mcl=8192$ at every $p\in\mathcal P_1$
(\cref{def:record:parameters}).

\paragraph{Support components (``pieces'') of a lift.}
Let $\mathcal P$ be a finite set of primes with $e_p=1$ for $p\in\mathcal P$, so
that $X_{\mathcal P}=\prod_{p\in\mathcal P}\Fp$, and put $q_U=\prod_{p\in U}p$
for $U\subseteq\mathcal P$; a character of $X_{\mathcal P}$ with exact support
$U$ has conductor $q_U$. For $f\colon I\to\C$ and $U\subseteq\mathcal P$, the
\emph{piece of exact support $U$} of the lift of $f$ on $X_{\mathcal P}$ is
$f_U(z)=\sum_\xi\widehat f_I(\xi)e(z\cdot\xi)$, the sum over the characters $\xi$
of $X_{\mathcal P}$ of exact support $U$. Each $f_U$
depends only on $z_U$ and has mean zero in every coordinate of $U$; if $f$ is
real, so is $f_U$. The mean of a lift is its coefficient at $\xi=0$,
$\E_z\mathcal L_{I,N}f(z)=\widehat f_I(0)=\E_If$, since $q(0)=1\le Q$. A lift of
a nonnegative function can take negative values: all positivity in this paper
comes from the reflection positivity of the laws, not from the inputs.

\subsection{Assumptions (R1)--(R5) on the per-atom laws (proved in Part~V)}\label{subsec:lifts:reference}

The analytic stage uses the following properties of the reference laws and of
the factors of the comparison law, and nothing else about the local laws.

\begin{hypothesis}[Reference atom laws]\label{hyp:record:reference}
For every interval scale $\Lsc\ge1$, with $N=\Kint\Lsc$:
\begin{enumerate}[label=(R\arabic*)]
\item\label{it:record:R1} every $\mu_\Blk$ is a probability measure on
 $X_\Blk^V$, and the law of an atom does not depend on $N$;
\item\label{it:record:R2} every $\mu_\Blk$ is reflection positive
 (\cref{def:rp}) at every cut of the family $\mathfrak F$, that is, at the
 chamber cuts of both chamber coordinates and at the reserved cut $\iota$; in
 particular at the seven cuts of the sorting sequence of \cref{lem:slot-array};
\item\label{it:record:R3} every $\mu_\Blk$ is square-affine invariant
 (\cref{def:map:invariance}): it is invariant under the simultaneous maps
 $z_{v,p}\mapsto r_pz_{v,p}+c_p$, where $r_p$ is a square unit and $c_p$
 an element of $\Z/p^{e_p}\Z$, chosen independently for $p\in\Blk$; and every
 signed measure $\rho_\Blk$ on $X_\Blk^V$ that occurs in Parts~III and~V at $\Blk$ as a component, a replacement, a one-layer folded copy, or a difference of such measures, has a variation $|\rho_\Blk|$ whose one-slot marginals are all
 $|\rho_\Blk|(1)\Unif(X_\Blk)$. Both hold for every measure that is invariant
 under these maps, since the variation of an invariant measure is invariant;
 in particular every one-slot marginal of $\mu_\Blk$ is $\Unif(X_\Blk)$;
\item\label{it:record:R4} for every tail prime $p\ge\Pfive$, with
 $\gamma=255/1024$: the law $\mu_p$ and its accepted part
 $\bar\nu_p=\nu_p/\nu_p(1)$, where $\nu_p$ is the factor at the one-prime atom $\{p\}$ of the comparison
 law $\nu$ of \cref{def:record:comparison-law} (at a tail prime the
 comparison law is nonnegative, \cref{lem:tail-prices}), are
 $(\gamma,v)$-mixing
 (\cref{def:mixing}) for every slot $v$; and for each query $(u,v)$ of the
 query family $\mathcal E_k$ of \cref{def:cp0}(c) (in Part~V, $\mathcal E_6$)
 the $(u,v)$-marginal of $\bar\nu_p$ is $\Unif(\Om)$ and
 $\bar\nu_p$ is $(\gamma,(u,v))$-mixing: its conditional operator from the
 pair $(u,v)$ to the other $70$ slots fixes constants, maps mean-zero functions
 to mean-zero functions and has norm at most $p^{-\gamma}$ on them;
\item\label{it:record:R5} under every accepted part, that is every normalized component $\bar\nu_{\Blk,i}=\nu_{\Blk,i}/\nu_{\Blk,i}(1)$ of positive mass of the comparison law (\cref{def:record:comparison-law}) (components of mass zero are omitted), almost surely each query pair is a
 strict-square pair at every prime of $\Blk$ with $e_p=1$, and at a clique
 prime $p\in\mathcal P_1$ its difference is nonzero, of even valuation
 $2j<e_p$, with a square unit part.
\end{enumerate}
\end{hypothesis}

Part~V verifies these properties: \ref{it:record:R1}--\ref{it:record:R3} for
the clique atoms by \cref{lem:frontier:iidlocal}, for the middle blocks by
\cref{lem:front12:reference}, for the upper blocks by
\cref{lem:upper-reference-rp} and for the tail atoms by \cref{lem:tail-law};
\ref{it:record:R4} is \cref{lem:record:tail-mixing}; \ref{it:record:R5} is the
acceptance property (a) of \cref{def:record:comparison-law}, checked range by
range and component by component; and the variation clause of \ref{it:record:R3} is part of the
admissibility of every selected law (\cref{def:admissible}). Each of the four
reference lemmas proves reflection positivity at every cut of $\mathfrak F$,
not only at the sorting cuts. This is needed: the one-cut arguments of
Parts~III and~V also use first-coordinate chamber cuts outside the sorting
sequence, for instance $\vartheta_{0,2}$, with plus half $\{2,3,4\}$, which is
neither a sorting cut nor a sorting cut with its halves exchanged.

\subsection{Averaging a lift over intervals and residues is exact}\label{subsec:lifts:index}

The external index replaces the lift by averages over subintervals
and residue classes. The next lemma says these averages return the
original lift exactly, with no error term; this is why conditioning
on the index preserves the mean $\alpha$.

\begin{lemma}[Index averages]\label{lem:twointerval}
Let $I$ be a nonempty interval of integers and $I=J_1\sqcup\dots\sqcup J_k$ a
partition of $I$ into nonempty intervals. For every $f\colon I\to\C$,
\[
 \mathcal L_{I,N}f=\sum_{j=1}^k\frac{|J_j|}{|I|}\,\mathcal L_{J_j,N}\big(f|_{J_j}\big).
\]
Consequently, for $\Lsc\ge1$, $N=\Kint\Lsc$, $A\subseteq[N]$ and a positive
divisor $D_0$ of $\Mzero$:
\begin{enumerate}[label=\textup{(\roman*)}]
\item\label{it:ti:intervals}
 $\Kint^{-1}\sum_{b\in\Z/\Kint\Z}g^{I_b,D_0,s}=g^{[N],D_0,s}$ for every
 $s\in\Z/D_0\Z$;
\item\label{it:ti:residues} if $s$ is uniform on $\Z/D_0\Z$ and $D'$ is a
 positive divisor of $D_0$, then
 $\E[g^{I,D_0,s}\mid s\bmod D']=g^{I,D',\,s\bmod D'}$ for every
 $I\in\mathcal I_\Lsc$; in particular ($D'=1$)
 $\E_sg^{I,D_0,s}=\mathcal L_{I,N}(\one_A|_I)$;
\item\label{it:ti:mean} if $b$ and $s$ are independent and uniform on
 $\Z/\Kint\Z$ and on $\Z/D_0\Z$, then $\E_{b,s}g^{I_b,D_0,s}=F_A$ and
 $\E_sg^{[N],D_0,s}=F_A$, identically on $G_N$.
\end{enumerate}
These identities are exact and hold for every fixed $\Kint\ge2$.
\end{lemma}

\begin{proof}
For every $\xi\in\Q/\Z$, $\widehat f_I(\xi)=\sum_j(|J_j|/|I|)
\widehat{(f|_{J_j})}_{J_j}(\xi)$, because the $J_j$ partition $I$. The lift sums
these coefficients over the set $\Xi_N$, which depends neither on the interval
nor on the function, so it is linear in the coefficient vector and the identity
follows. \ref{it:ti:intervals} Apply it to the partition of $[N]$ into the
$I_b$, for which $|I_b|/N=1/\Kint$, and $f=D_0\one_{n\equiv s}\one_A$.
\ref{it:ti:residues} Let $t\in\Z/D'\Z$ and $n\in\Z$. Given $s\equiv t\pmod{D'}$,
$s$ is uniform among the $D_0/D'$ residues modulo $D_0$ above $t$, of which
exactly one is $\equiv n\pmod{D_0}$ if $n\equiv t\pmod{D'}$, and none otherwise;
hence the conditional expectation of $D_0\one_{n\equiv s\,(D_0)}$ is
$D'\one_{n\equiv t\,(D')}$. Multiply by $\one_A(n)$, restrict to $I$ and apply
the linear map $\mathcal L_{I,N}$. \ref{it:ti:mean} By \ref{it:ti:residues},
$\E_sg^{I,D_0,s}=\mathcal L_{I,N}(\one_A|_I)$, which is $F_A$ for $I=[N]$; averaging
over $b$ and using the partition identity gives
$\E_{b,s}g^{I_b,D_0,s}=\Kint^{-1}\sum_b\mathcal L_{I_b,N}(\one_A|_{I_b})=F_A$.
\end{proof}

Only the linearity of the lift and the independence of $\Xi_N$ from the
interval and the input are used; the prime-power coordinates of $G_N$ play no
role. In the index law (\cref{def:hyb:coupled}) the residue $s\in\Z/\Mzero\Z$
has prime-power coordinates by the Chinese remainder theorem. Revealing a
residue to depth $J_p$ at each prime is conditioning on
$s\bmod D_{\boldsymbol J}$, where $D_{\boldsymbol J}=\prod_p p^{J_p}$.
When every revealed depth $J_p$ is either full (the whole
$p$-coordinate of $\Mzero$) or zero, write $W$ for the set of fully
revealed coordinates and $D_W$ for the corresponding modulus; in either case
\ref{it:ti:residues} turns $g^{I,\Mzero,s}$ into $g^{I,D_{\boldsymbol J},s\bmod D_{\boldsymbol J}}$. For
inputs with independent indices at distinct slots, such as the pairwise
disjoint endpoints of the queries, the expectation and the conditional
expectations of the product are the products of those of the factors.

\subsection{\texorpdfstring{$L^{2r}$}{L\textasciicircum (2r)} bounds for lifts, linear in the density}\label{subsec:lifts:moments}

The moment bounds have three ingredients: a conditional Parseval bound for the
pieces of a lift, a recursion for their square function, and decoupling over
the prime coordinates (compare the general decoupling formalism for
multilinear and U-statistic averages in \cite{dlPMS}; the coordinate-colour
argument needed here is proved directly below). We state them separately, so
that the dependence on the conductor cutoff and on the number of prime factors
remains visible. The first ingredient is arithmetic. A generic hypercontractive
bound would not do: on $\Fp$ the mean-zero function $\one_{x=0}-1/p$ has
$\|f\|_{2r}/\|f\|_2$ of order $p^{1/2-1/(2r)}$. The pieces of a lift come from an
interval, and the Gram (large-sieve) bound uses this.

\begin{lemma}[Gram bound and conditional square sums]
\label{lem:lift:conditional-squares}
Let $I$ be an interval of integers and $Q\ge1$.
\begin{enumerate}[label=(\alph*)]
\item If $\Xi$ is a set of distinct elements of $\Q/\Z$ with reduced
 denominators at most $Q$, then for every $f\colon I\to\C$
 \[
  \sum_{\xi\in\Xi}|\widehat f_I(\xi)|^2\le\Big(1+\frac{Q^4}{2|I|}\Big)\E_I|f|^2 .
 \]
\item Let $|I|\ge10Q^6$ and $\|f\|_\infty\le M$, let $\mathcal P$ be a finite
 set of primes with $e_p=1$, and let $\mathcal F$ be a family of subsets
 $U\subseteq\mathcal P$ with $q_U\le Q$. For every $W\subseteq\mathcal P$ with
 $q_W\le Q$ and every $z_W\in X_W$, the pieces $f_U$ of the lift of $f$ on
 $X_{\mathcal P}$ satisfy
 \begin{equation}\label{eq:cond-square}
  \sum_{U\in\mathcal F,\,U\supseteq W}\E\big[|f_U|^2\ \big|\ z_W\big]
  \le4M^2\,4^{|W|}.
 \end{equation}
\end{enumerate}
\end{lemma}

\begin{proof}
(a) Distinct elements of $\Q/\Z$ with reduced denominators at most $Q$ are at
distance at least $Q^{-2}$, and there are at most $\sum_{q\le Q}\varphi(q)\le Q^2$
of them. Write $\|t\|$ for the distance from $t\in\Q/\Z$ to $0$.
For $\xi\ne\xi'$ in $\Xi$,
$|I|^{-1}|\sum_{n\in I}e(n(\xi-\xi'))|\le(2|I|\,\|\xi-\xi'\|)^{-1}
\le Q^2/(2|I|)$. So the Gram matrix $G$ of the unit vectors
$\phi_\xi=(|I|^{-1/2}e(n\xi))_{n\in I}$, $\xi\in\Xi$, has norm at most
$1+Q^4/(2|I|)$ (Schur's test), and Bessel's inequality in Gram form,
$\sum_\xi|\ip f{\phi_\xi}|^2\le\|G\|\,\|f\|^2$, is the claim, since
$|\ip f{\phi_\xi}|^2=|I|\,|\widehat f_I(\xi)|^2$ and $\|f\|^2=|I|\E_I|f|^2$.

(b) Write $U=W\sqcup U'$ and $\xi=\xi_W+\eta$ with exact supports $W$ and $U'$.
After fixing $z_W$, the coefficient of $e(z_{U'}\cdot\eta)$ in $f_U$ is
\[
 c_W(\eta;z_W)=\frac1{|I|}\sum_{n\in I}f(n)e(-n\eta)\prod_{p\in W}
 \big(p\one_{n\equiv z_p\,(p)}-1\big),
\]
since $\sum_{a=1}^{p-1}e(a(z_p-n)/p)=p\one_{n\equiv z_p}-1$. By conditional
Parseval the left side of \eqref{eq:cond-square} is
$\sum_{\eta\in\Omega_W}|c_W(\eta;z_W)|^2$, summed over the union $\Omega_W$ of the
sets of characters of exact support $U\setminus W$, $U\in\mathcal F$,
$U\supseteq W$ (different $U$ give disjoint sets); the elements of $\Omega_W$
have reduced denominators at most $Q/q_W$, all coprime to $q_W$. Expanding the
product over the subsets $J\subseteq W$ gives
$c_W=\sum_J(-1)^{|W\setminus J|}c_J$ with
$c_J(\eta)=(q_J/|I|)\sum_{n\in I,\,n\equiv z_J\,(q_J)}f(n)e(-n\eta)$. Write
$n=n_J+q_Jn'$ with $n'$ in an interval $T_J$ of length $L_J$, where
$|L_J-|I|/q_J|\le1$; since $|I|/q_J\ge10Q^5$, $L_J\ge0.9|I|/q_J$ and
$q_JL_J\le1.1|I|$. The frequencies $q_J\eta$, $\eta\in\Omega_W$, are distinct,
because $q_J$ is a unit modulo every denominator in $\Omega_W$, and have reduced
denominators at most $Q$. With $F_J(n')=f(n_J+q_Jn')$,
$|c_J(\eta)|=(q_JL_J/|I|)\,|\widehat{(F_J)}_{T_J}(q_J\eta)|$, and (a) on $T_J$, with
$1+Q^4/(2L_J)\le1+Q^5/(1.8|I|)\le19/18$, gives
$(\sum_\eta|c_J(\eta)|^2)^{1/2}\le1.1\sqrt{19/18}\,M\le2M$. The triangle inequality
over the $2^{|W|}$ sets $J$ gives
$(\sum_\eta|c_W(\eta;z_W)|^2)^{1/2}\le2^{|W|}\cdot2M$, which is
\eqref{eq:cond-square}.
\end{proof}

\begin{lemma}[Square-function recursion]\label{lem:sqfun-rec}
Let $X=\prod_{p\in\mathcal P}X_p$ be a finite product of finite probability
spaces, $\mathcal F$ a finite family of subsets of $\mathcal P$ of size at most
$k$, and $h_U\colon X\to\C$ $(U\in\mathcal F)$ functions such that $h_U$
depends only on $z_U$. Suppose that for some $c\ge0$, every $W\subseteq\mathcal P$
contained in a member of $\mathcal F$, and every $z_W$,
\[
 \sum_{U\in\mathcal F,\,U\supseteq W}\E\big[|h_U|^2\ \big|\ z_W\big]\le c\,4^{|W|}.
\]
Put $S=(\sum_{U\in\mathcal F}|h_U|^2)^{1/2}$ and $T_j=\E S^{2j}$, so $T_0=1$.
Then $T_j\le c\,4^k2^{(j-1)k}\,T_{j-1}$ for every $j\ge1$; hence, for all
integers $0\le j_0\le r$,
\[
 T_r\le(c\,4^k)^{r-j_0}\,2^{k(r(r-1)-j_0(j_0-1))/2}\,T_{j_0}.
\]
\end{lemma}

\begin{proof}
Let $j\ge1$. Expand $S^{2j}=\sum_{U_1,\dots,U_j}\prod_i|h_{U_i}|^2$ and fix
$U_2,\dots,U_j$, with union $\mathcal U$, so $|\mathcal U|\le(j-1)k$. The
product $\prod_{i\ge2}|h_{U_i}|^2$ depends only on $z_{\mathcal U}$, and since the
coordinates are independent,
$\E[|h_{U_1}|^2\mid z_{\mathcal U}]=\E[|h_{U_1}|^2\mid z_W]$ with
$W=U_1\cap\mathcal U$. Grouping the supports $U_1$ by $W$, which takes at most
$2^{|\mathcal U|}\le2^{(j-1)k}$ values, each contained in a member of
$\mathcal F$ and of size at most $k$, gives pointwise
$\sum_{U_1}\E[|h_{U_1}|^2\mid z_{\mathcal U}]\le2^{(j-1)k}c\,4^k$. Multiplying by
$\prod_{i\ge2}|h_{U_i}|^2$, taking expectations and summing over
$U_2,\dots,U_j$ gives $T_j\le c4^k2^{(j-1)k}T_{j-1}$. Iterate from $j_0+1$ to
$r$, using $\sum_{j=j_0+1}^r(j-1)=(r(r-1)-j_0(j_0-1))/2$.
\end{proof}

For the pieces of a lift with $\|f\|_\infty\le M$, \eqref{eq:cond-square} gives
the hypothesis with $c=4M^2$, and the case $j_0=0$ reads
$\|S\|_{2r}\le2M\,2^{k(r+3)/4}$.

\begin{lemma}[Decoupling by coordinate colours]\label{lem:lift:decoupling}
Let $X=\prod_{p\in\mathcal P}X_p$ be a finite product of finite probability
spaces, $\mathcal F$ a finite family of subsets of $\mathcal P$ of size at most
$k$, and $f_U\colon X\to\C$ $(U\in\mathcal F)$ functions such that $f_U$
depends only on the coordinates in $U$ and has mean zero in each of them. Put
$S=(\sum_{U\in\mathcal F}|f_U|^2)^{1/2}$ and, for an integer $r\ge1$,
$K_{2r}=((2r)!/(2^rr!))^{1/(2r)}$. Then
\[
 \Big\|\sum_{U\in\mathcal F}f_U\Big\|_{2r}\le(k+1)(2eK_{2r})^k\|S\|_{2r}.
\]
\end{lemma}

\begin{proof}
We use Khintchine's inequality $\|\sum a_i\epsilon_i\|_{2r}\le
K_{2r}(\sum|a_i|^2)^{1/2}$ for independent uniform signs $\epsilon_i$. For real
$a_i$ the even moment is at most that of $\sum a_i\mathrm g_i$ with independent
standard Gaussians (expand: only even exponents survive, with nonnegative
coefficients, and $\E\epsilon^{2m}=1\le\E\mathrm g^{2m}$), which is
$(2r-1)!!(\sum a_i^2)^r=K_{2r}^{2r}(\sum a_i^2)^r$. For complex $a_i$ apply this
to the real and imaginary parts and Minkowski in $L^r$ to
$|X|^2=(\mathrm{Re}X)^2+(\mathrm{Im}X)^2$.

The same constant applies to Hilbert-space vectors $v_i$. Work in their
finite-dimensional real span (realifying a complex Hilbert space), and let
$G$ be a standard Gaussian vector there. With $c_r=(2r-1)!!=K_{2r}^{2r}$,
the scalar inequality and then Minkowski in $L^r$ give
\[
 \E_\epsilon\Big\|\sum_i\epsilon_i v_i\Big\|^{2r}
 =c_r^{-1}\E_G\E_\epsilon
       \Big\langle G,\sum_i\epsilon_i v_i\Big\rangle^{2r}
 \le\E_G\Big(\sum_i\langle G,v_i\rangle^2\Big)^r
 \le c_r\Big(\sum_i\|v_i\|^2\Big)^r.
\]
This also holds conditionally when the vectors depend on other variables.

Fix $j\ge1$ and colour the coordinates independently and uniformly with $j$
colours; a set of size $j$ is \emph{transversal} if its elements have distinct
colours, which happens with probability $\pi_j=j!/j^j$. Fix a colouring and let
$\mathcal T$ be the transversal members of $\mathcal F$ of size $j$. Take
independent copies $X^0,X^1$ of the coordinates and put
$D_U=\prod_{p\in U}(I-\mathrm{swap}_p)f_U(X^0)$, where $\mathrm{swap}_p$
exchanges the $p$-coordinates of $X^0$ and $X^1$. Then
$D_U=\sum_{U'\subseteq U}(-1)^{|U'|}f_U(X^{U'})$, where $X^{U'}$ takes the
coordinates in $U'$ from $X^1$ and the others from $X^0$; integrating out $X^1$
makes every term with $U'\ne\emptyset$ vanish, since $f_U$ has mean zero in
each coordinate of $U$. So
$\E_{X^1}D_U=f_U(X^0)$, and by Jensen
$\|\sum_{\mathcal T}f_U\|_{2r}\le\|\sum_{\mathcal T}D_U\|_{2r}$. Each
$\mathrm{swap}_p$ preserves the joint law of $(X^0,X^1)$ and changes the sign of
$D_U$ exactly when $p\in U$; hence
$\|\sum_{\mathcal T}D_U\|_{2r}=\|\sum_{\mathcal T}\epsilon_UD_U\|_{2r}$, where
$\epsilon_U=\prod_{p\in U}\epsilon_p$ for independent uniform signs $\epsilon_p$.
A transversal $U$ has exactly one coordinate $p_c(U)$ of each colour $c$, so
$\epsilon_U=\prod_{c=1}^j\epsilon_{p_c(U)}$, and we integrate the signs one colour
at a time. With the signs of colour $1$ free and the rest fixed,
$\sum_{\mathcal T}\epsilon_UD_U=\sum_p\epsilon_pA_p$ over the coordinates of colour
$1$, and Khintchine gives the bound $K_{2r}\|\sum_p|A_p|^2\|_r^{1/2}$. Each $A_p$
is in turn $\sum_{p'}\epsilon_{p'}A_{p,p'}$ over the coordinates of colour $2$,
and the Hilbert-space inequality applied to the vectors
$(A_{p,p'})_p$, conditionally on all remaining variables, gives
$\|\sum_p|A_p|^2\|_r\le K_{2r}^2\|\sum_{p,p'}|A_{p,p'}|^2\|_r$. After the $j$
colours,
\[
 \Big\|\sum_{\mathcal T}\epsilon_UD_U\Big\|_{2r}
 \le K_{2r}^j\Big\|\Big(\sum_{U\in\mathcal T}|D_U|^2\Big)^{1/2}\Big\|_{2r}.
\]
For transversal $U$ the sets $U'\subseteq U$ correspond to sets $C$ of colours,
so $D_U=\sum_{C\subseteq[j]}(-1)^{|C|}f_U(X^{(C)})$, where $X^{(C)}$ takes the
coordinates of the colours in $C$ from $X^1$ and the others from $X^0$ and does
not depend on $U$. As every $X^{(C)}$ has the law of $X^0$, Minkowski in
$\ell^2(\mathcal T)$ and in $L^{2r}$ gives
$\|(\sum_{\mathcal T}|D_U|^2)^{1/2}\|_{2r}\le2^j\|S\|_{2r}$. So
$\|\sum_{\mathcal T}f_U\|_{2r}\le(2K_{2r})^j\|S\|_{2r}$ for every colouring. Every
$U$ of size $j$ is transversal with probability $\pi_j$, so
$\sum_{|U|=j}f_U=\pi_j^{-1}\E_{\rm col}\sum_{\mathcal T}f_U$, and Minkowski over
the colourings gives
$\|\sum_{|U|=j}f_U\|_{2r}\le\pi_j^{-1}(2K_{2r})^j\|S\|_{2r}\le(2eK_{2r})^j\|S\|_{2r}$,
as $j^j/j!\le e^j$. The term with $|U|=0$ is constant and bounded in modulus by
$S$. Sum over $0\le j\le k$, using $2eK_{2r}\ge1$.
\end{proof}

The moment theorem combines the three lemmas. It keeps the $\ell^2$ mass
$\E_I|f|^2$ of the integer input, and with it the density.

\begin{lemma}[Moments that retain the $\ell^2$ mass]
\label{lem:record:lift-moments}
Let $\mathcal S'\subseteq\mathcal P_1$ and let $\mathcal P$ be a finite set of
primes with $e_p=1$; put $X=X_{\mathcal S'}\times X_{\mathcal P}$. Let $Q\ge1$, let
$I$ be an interval with $|I|\ge10Q^6$, and let $f\colon I\to\C$ with
$\|f\|_\infty\le M$. For each character $\chi$ of $X_{\mathcal S'}$ with
$q(\chi)\le Q$ let $\mathcal F_\chi$ be a family of sets $U\subseteq\mathcal P$
with $q(\chi)q_U\le Q$, and put
\[
 G(z',z)=\sum_\chi e(z'\cdot\chi)\sum_{U\in\mathcal F_\chi}(f_\chi)_U(z),
 \qquad f_\chi(n)=f(n)e(-n\chi),
\]
where $(f_\chi)_U$ is the piece of exact support $U$ of the lift of $f_\chi$ on
$X_{\mathcal P}$. We call such a $G$ a \emph{whole-support piece} of the lift of
$f$ on $X$ with cutoff $Q$. Let $k=\max\{|U|:U\subseteq\mathcal P,\ q_U\le Q\}$,
and let $K_{2r}$ be as in \cref{lem:lift:decoupling}.
\begin{enumerate}[label=(\alph*)]
\item $\|G\|_2^2\le(1+Q^4/(2|I|))\,\E_I|f|^2\le\tfrac{21}{20}\E_I|f|^2$.
\item For every integer $r\ge1$,
 \[
  \|G\|_{2r}\le D_{\mathcal S'}(k+1)(2eK_{2r})^k(2M)^{1-1/r}
  2^{k(r^2+3r-4)/(4r)}\big(\tfrac{21}{20}\E_I|f|^2\big)^{1/(2r)}.
 \]
\item For fixed $r$ and $\epsilon>0$ (an auxiliary exponent, not the exponent
 of the main theorem) there is $C_{r,\epsilon}$, depending only
 on $r$, $\epsilon$ and $Q_{\rm cl}$, with
 $\|G\|_{2r}^{2r}\le C_{r,\epsilon}Q^\epsilon M^{2r-2}\E_I|f|^2$. In particular,
 for $f=D_0\one_{n\equiv s\,(D_0)}\one_A$ on $I$ with $D_0\mid\Mzero$, the
 right-hand side is at most $C_{r,\epsilon}Q^\epsilon\Mzero^{2r}\alpha_I$, where
 $\alpha_I=|A\cap I|/|I|$.
\end{enumerate}
No division by $\E_I|f|^2$, $M$ or a density occurs; if $\E_I|f|^2=0$ then
$G=0$.
\end{lemma}

The lift itself is such a piece: a character of $G_N$ is $\chi+\xi$ with $\chi$
a character of $X_{\mathcal P_1}$ and $\xi$ one of $X_{\mathcal P_3\cup\mathcal T_N}$,
$q(\chi+\xi)=q(\chi)q(\xi)$ and $\widehat f_I(\chi+\xi)=\widehat{(f_\chi)}_I(\xi)$;
so $\mathcal L_{I,N}f$ is the whole-support piece with $\mathcal S'=\mathcal P_1$,
$\mathcal P=\mathcal P_3\cup\mathcal T_N$ and $\mathcal F_\chi=\{U:q(\chi)q_U\le Q\}$.

\begin{proof}
\emph{Step 1: the Gram bound.} The Fourier coefficient of $G$ at $\chi+\xi''$ is
$0$ or $\widehat f_I(\chi+\xi'')$, and the characters with a nonzero coefficient
are distinct elements of $\Q/\Z$ with reduced denominators at most $Q$; so
Parseval and \cref{lem:lift:conditional-squares}(a) prove (a), as
$1+Q^4/(2|I|)\le21/20$. Applied to $f_\chi$, whose modulus is that of $f$, with
the cutoff $Q/q(\chi)$, the same argument gives for each $\chi$
\begin{equation}\label{eq:record:Echi}
 E_\chi:=\sum_{U\in\mathcal F_\chi}\|(f_\chi)_U\|_2^2\le\tfrac{21}{20}\E_I|f|^2 .
\end{equation}

\emph{Step 2: one prefix character.} Fix $\chi$ and write $h_U=(f_\chi)_U$,
$S^2=\sum_{U\in\mathcal F_\chi}|h_U|^2$ and $T_j=\E S^{2j}$, so $T_1=E_\chi$. By
\eqref{eq:cond-square} for $f_\chi$, with $\|f_\chi\|_\infty\le M$ and
$|I|\ge10Q^6$, the hypothesis of \cref{lem:sqfun-rec} holds with $c=4M^2$; the
case $j_0=1$ gives $T_r\le(2M)^{2r-2}2^{k(r^2+3r-4)/2}E_\chi$, because
$(4M^24^k)^{r-1}2^{kr(r-1)/2}=(2M)^{2r-2}2^{k(r-1)(r+4)/2}$. The $h_U$ depend only
on $z_U$ and have mean zero in each coordinate of $U$, so
\cref{lem:lift:decoupling} gives
\[
 \Big\|\sum_Uh_U\Big\|_{2r}\le(k+1)(2eK_{2r})^k\,T_r^{1/(2r)}
 \le(k+1)(2eK_{2r})^k(2M)^{1-1/r}2^{k(r^2+3r-4)/(4r)}E_\chi^{1/(2r)} .
\]

\emph{Step 3: the prefix.} $|e(z'\cdot\chi)|=1$, and there are at most
$D_{\mathcal S'}$ characters $\chi$. Minkowski over $\chi$, Step~2 and
\eqref{eq:record:Echi} prove (b).

\emph{Step 4: asymptotics.} A product of $k$ distinct primes is at least
$(k+1)!$, so $(k+1)!\le Q$ and $k=O(\log Q/\log\log Q)$. Hence every factor of
(b) that is exponential in $k$ is $Q^{o(1)}$ for fixed $r$, and
$D_{\mathcal S'}\le Q_{\rm cl}$. Raising (b) to the power $2r$ gives the first
part of (c). For $f=D_0\one_{n\equiv s}\one_A$ we have $M=D_0\le\Mzero$ and
$\E_I|f|^2=D_0^2\#\{n\in A\cap I:n\equiv s\}/|I|\le\Mzero^2\alpha_I$.
\end{proof}

In \cref{sec:scale-comparison} the lemma is applied to the inputs on $G_N$,
whose specializations at fixed labels on an arbitrary set of atoms are integrated over the
labels by Fubini, $\E_{z_{\mathcal W}}\|G(z_{\mathcal W},\cdot)\|_s^s=\|G\|_s^s$
for every set $\mathcal W$ of coordinates and every $s\ge1$ (the fibres of a
finite uniform product have equal mass); and to the lifts of rescaled sets on
the complement of the selected atoms (\cref{lem:record:fibre-scale}(e)). A fixed
label, without averaging, satisfies the explicit bound of the next
corollary.

\begin{corollary}[Conditional moments]\label{lem:record:conditional-moments}
In the setting of \cref{lem:record:lift-moments}, let
$\mathcal W\subseteq\mathcal P$, let $z_{\mathcal W}\in X_{\mathcal W}$, and fix arbitrary values on any subset of the coordinates of $X_{\mathcal S'}$; the bound below holds for every such choice. Then for every
integer $r\ge1$, with $A_r=2eK_{2r}2^{(r+3)/4}$,
\[
 \|G(z_{\mathcal W},\cdot)\|_{2r}\le D_{\mathcal S'}\,2M\,3^{|\mathcal W|}(k+1)A_r^k,
\]
the norm being taken over the remaining coordinates with their uniform law. If
$Q\le N^\beta$ and $q_{\mathcal W}\le N^{C_0}$ for a fixed $C_0$, the right-hand
side is at most $C_{r,\epsilon}MN^\epsilon$ for every $\epsilon>0$, uniformly in
the labels.
\end{corollary}

\begin{proof}
Fixing prefix coordinates multiplies the term of $\chi$ by a unimodular
constant or leaves it a function of the free prefix coordinates, so by
Minkowski over $\chi$ it suffices to treat one $\chi$, with $h_U=(f_\chi)_U$.
Group the supports by $U_0=U\cap\mathcal W$. For fixed $U_0$ the functions
$h'_{U'}(y)=h_{U_0\cup U'}(z_{U_0},y)$, $U'\subseteq\mathcal P\setminus\mathcal W$,
depend only on $y_{U'}$ and have mean zero in each coordinate of $U'$; for
$W'\subseteq\mathcal P\setminus\mathcal W$ their conditional square sum over
$U'\supseteq W'$ given $y_{W'}$ is part of the sum in \eqref{eq:cond-square} with
pinned set $U_0\cup W'$, hence at most $4M^24^{|U_0|}4^{|W'|}$.
\Cref{lem:sqfun-rec} with $c=4M^24^{|U_0|}$ and $j_0=0$, and
\cref{lem:lift:decoupling}, give
$\|\sum_{U'}h'_{U'}\|_{2r}\le2M2^{|U_0|}(k+1)A_r^k$. Summing over
$U_0\subseteq\mathcal W$, with $\sum_{U_0}2^{|U_0|}=3^{|\mathcal W|}$, proves the
bound. A product of $|\mathcal W|$ distinct primes is at least
$(|\mathcal W|+1)!$, so $3^{|\mathcal W|}$ and the factors exponential in $k$ are
$N^{o(1)}$.
\end{proof}

\begin{remark}[Averages and suprema over fibres]\label{rem:lift:fibres}
The estimates of \cref{sec:scale-comparison} use only averaged moments.
Averaging the fibre norms over $z_{\mathcal W}$ returns the full norm exactly.
For a general function a single fixed fibre can be larger by the factor
$q_{\mathcal W}^{1/s}$, the reciprocal of the $s$-th root of the fibre probability $q_{\mathcal W}^{-1}$:
$\|G(z_{\mathcal W},\cdot)\|_s\le q_{\mathcal W}^{1/s}\|G\|_s$; for a lift a
fixed label costs only $3^{|\mathcal W|}=N^{o(1)}$. This does not bound a pinned
lift by a pinned density, and a pinned lift need not vanish on an empty fibre
(\cref{rem:record:fibres}).
\end{remark}

\subsection{Tensor products preserve exact Fourier support}\label{subsec:lifts:sectors}

A tensor product of operators that fix the constants on both sides preserves
the orthogonal decomposition of a product space by exact nonconstant support,
so its norm, the norm of a difference of two such products, and their symmetric
and skew parts can be computed support by support.

\paragraph{Setting.} Let $\mathcal K$ be a finite set. For
$j\in\{0\}\cup\mathcal K$ let $(\Omega_j,\pi_j)$ and $(\Omega'_j,\pi'_j)$ be
finite probability spaces, $L^2(\pi_j)$ the (real or complex) functions on the
support of $\pi_j$, and $L^2_0(\pi_j)=\{f:\int f\,d\pi_j=0\}$; put
$\Omega=\prod_j\Omega_j$, $\pi=\bigotimes_j\pi_j$ and
$\pi'=\bigotimes_j\pi'_j$. The index $0$ is a
\emph{spectator}, possibly a one-point space. For $S\subseteq\mathcal K$ the
\emph{sector} of exact support $S$ is
\[
 H_S(\pi)=L^2(\pi_0)\otimes\bigotimes_{j\in S}L^2_0(\pi_j)\otimes
 \bigotimes_{j\in\mathcal K\setminus S}\C\one ,
\]
and $X_S$ is the component of $X\in L^2(\pi)$ in $H_S(\pi)$, taken
coordinatewise for maps into a finite-dimensional space. The sectors are
mutually orthogonal and sum to $L^2(\pi)$, and for $W\subseteq\mathcal K$ the
conditional expectation $\E_\pi[X\mid\omega_0,\omega_W]$ is the orthogonal
projection onto $\bigoplus_{S\subseteq W}H_S(\pi)$, that is, it equals
$\sum_{S\subseteq W}X_S$. A linear operator $T_j\colon L^2(\pi_j)\to L^2(\pi'_j)$
is \emph{sector-preserving} if $T_j\one=\one$ and $T_j^*\one=\one$. Then $T_j$
maps $L^2_0(\pi_j)$ into $L^2_0(\pi'_j)$, as
$\ip{\one}{T_jf}=\ip{T_j^*\one}f=\ip{\one}f$; we write $T_j^\circ$ for this
restriction and $\delta_j=\|T_j^\circ\|$ (this local $\delta_j$
is a norm, not one of the depth tails $\delta_{p,j}$). The Markov
operator
$T_jf(x')=\E[f(x)\mid x']$ of a coupling of $\pi'_j$ and $\pi_j$ (a probability
measure on $\Omega'_j\times\Omega_j$ with these marginals) is sector-preserving,
with $\delta_j\le1$, its maximal correlation; so is a doubly stochastic matrix
on a finite set with the uniform law. The spectator operator
$T_0\colon L^2(\pi_0)\to L^2(\pi'_0)$ is any linear operator.

\begin{lemma}[Sector lemma]\label{lem:sector}
In this setting let $T_j$ $(j\in\mathcal K)$ be sector-preserving and
$T=T_0\otimes\bigotimes_{j\in\mathcal K}T_j$; put $\delta_S=\prod_{j\in S}\delta_j$.
\begin{enumerate}[label=(\alph*)]
\item\emph{Sectors.} $T$ maps $H_S(\pi)$ into $H_S(\pi')$ and equals
 $T_0\otimes\bigotimes_{j\in S}T_j^\circ\otimes\bigotimes_{j\notin S}1$ there,
 of norm $\|T_0\|\,\delta_S$; the same holds for $T^*$. Consequently, if $R$ is
 a linear combination of such tensor products and their adjoints, all acting
 between the same two spaces, then $R$ maps every sector into the
 corresponding sector and $\|R\|=\max_S\|R|_{H_S(\pi)}\|$; when $\pi'=\pi$ this
 applies to the symmetric and skew parts $\frac12(R\pm R^*)$, and an
 $n\times n$ operator matrix whose entries are such combinations maps every
 $H_S(\pi)\otimes\C^n$ into itself; it is positive semidefinite if and only if
 each of these restrictions is.
\item\emph{Differences.} Let $T'_j$ $(j\in\mathcal K)$ be sector-preserving
 as well, with $\delta'_j=\|T'^\circ_j\|$, $\delta'_S=\prod_{j\in S}\delta'_j$,
 $\|T'^\circ_j-T^\circ_j\|\le d_j$ and $\max(\delta_j,\delta'_j)\le\kappa_j$, and
 $T'=T_0\otimes\bigotimes_jT'_j$.
 Then $T'-T$ vanishes on $H_\emptyset(\pi)$ and, on $H_S(\pi)$,
 \[
  \|(T'-T)|_{H_S}\|\le\|T_0\|\min\Big(\sum_{j\in S}d_j\!\!\prod_{k\in S\setminus\{j\}}\!\!\kappa_k,\
  \delta'_S+\delta_S\Big).
 \]
 Let moreover $\pi'=\pi$ and $T_0=I$. On $H_{\{j\}}$ the difference and its
 symmetric and skew parts are those of $T'^\circ_j-T^\circ_j$. If every $T_j$
 is self-adjoint, the skew part of $T'-T$ on $H_S$ is that of $T'|_{H_S}$, of
 norm at most $\delta'_S$. If $d_j\le d$ and $\kappa_j\le\kappa\le1/2$ for all
 $j$, the first bound is at most $|S|d\kappa^{|S|-1}\le2d\kappa\le d$ for
 $|S|\ge2$; so if also $\|\mathrm{Sym}(T'_j-T_j)\|\le d_{\rm s}$ and
 $\|\mathrm{Skew}(T'_j-T_j)\|\le d_{\rm a}$ for all $j$, then
 \[
  \|\mathrm{Sym}(T'-T)\|\le\max(d_{\rm s},2d\kappa),\qquad
  \|\mathrm{Skew}(T'-T)\|\le\max(d_{\rm a},2d\kappa),\qquad\|T'-T\|\le d .
 \]
 More generally, with $\pi'=\pi$ and $T_0=I$, the norms of $T'-T$ and of its
 symmetric and skew parts are the maxima over $S\ne\emptyset$ of the norms of
 their restrictions to $H_S$; on $H_{\{j\}}$ these restrictions are the
 corresponding parts of $T'^\circ_j-T^\circ_j$, and on $H_S$ with $|S|\ge2$
 their norms are at most $\delta_S+\delta'_S$, which is at most
 $\delta_j\delta_k+\delta'_j\delta'_k$ for any $j\ne k$ in $S$ when all
 $\delta_i,\delta'_i\le1$. Hence, in that case,
 \[
  \|\mathrm{Sym}(T'-T)\|\le\max\Bigl(\max_j\|\mathrm{Sym}(T'^\circ_j-T^\circ_j)\|,\
  \max_{j\ne k}(\delta_j\delta_k+\delta'_j\delta'_k)\Bigr),
 \]
 the second maximum being $0$ if $|\mathcal K|=1$, and the same holds for
 $\mathrm{Skew}(T'-T)$ and for $T'-T$.
\item\emph{Markov tensor maximum.} Let the spectator be trivial,
 $\mathcal K\ne\emptyset$, and $\delta_j\le1$ for all $j$ (for instance, Markov
 operators of couplings with exact marginals). Then
 $\|T-\Jproj\|=\max_j\delta_j$, where $\Jproj f=(\int f\,d\pi)\one$. If moreover
 $\pi'=\pi$ and each $T_j$ is self-adjoint with spectrum in $[-c_j,1]$,
 $0\le c_j\le1$, then $T\succeq-(\max_jc_j)I$.
\item\emph{Random inclusion.} Let $\Jc\subseteq\mathcal K$ be a random set
 containing each $j$ independently with probability $\delta_j\in[0,1]$, where
 here the $\delta_j$ are any numbers (in the applications, bounds for the norms
 of (a)), and $\delta_S=\prod_{j\in S}\delta_j$. Let $E$ be
 a finite-dimensional space with a positive semidefinite Hermitian form $M$,
 and $\|Y\|_M^2=\int M(Y(\omega),Y(\omega))\,d\pi(\omega)$ for $Y\colon\Omega\to E$.
 Then components of different supports are orthogonal for this form, and for
 every $X\colon\Omega\to E$
 \[
  \E_\Jc\big\|\E_\pi[X\mid\omega_0,\omega_\Jc]\big\|_M^2=\sum_{S\subseteq\mathcal K}\delta_S\|X_S\|_M^2 .
 \]
\item\emph{Exact-support decay.} For every $X\in L^2(\pi)$,
 $\|TX\|^2=\sum_S\|TX_S\|^2\le\|T_0\|^2\sum_S\delta_S^2\|X_S\|^2$.
\item\emph{Product densities.} If the spectator is trivial and
 $\varphi_j\in L^2(\pi_j)$ with $\int\varphi_j\,d\pi_j=1$ $(j\in\mathcal K)$, then
 $\|\bigotimes_j\varphi_j-\one\|^2=\prod_j(1+\|\varphi_j-\one\|^2)-1$.
\item\emph{Tensor spectrum} (no sector hypothesis). If $A_1,\dots,A_r$ are
 self-adjoint operators on finite-dimensional Hilbert spaces with spectra in
 $[-c_i,1+b_i]$, where $b_i\ge0$ and $0\le c_i\le1+b_i$, then the spectrum of
 $\bigotimes_iA_i$ lies in
 \[
  \Big[-\max_i\Big(c_i\prod_{k\ne i}(1+b_k)\Big),\ \prod_i(1+b_i)\Big];
 \]
 in particular
 $\bigotimes_iA_i\succeq-(\max_ic_i)\prod_i(1+b_i)I\succeq-(\max_ic_i)e^{\sum_ib_i}I$.
\item\emph{Two supports.} Let $\pi'=\pi$, $T_0=I$, let $T_j$ and $T'_j$ be as
 in \textup{(b)} with every $T_j$ self-adjoint, and write $T'_j=T_j+S_j+A_j$
 with $S_j=\mathrm{Sym}(T'_j-T_j)$ and $A_j=\mathrm{Skew}(T'_j-T_j)$, so that
 $T_j+S_j=\mathrm{Sym}\,T'_j$ and $A_j=\mathrm{Skew}\,T'_j$. Then $S_j$ and
 $A_j$ vanish on the constants and map $L^2_0(\pi_j)$ into itself, and on
 $H_{\{j,k\}}$, $j\ne k$, with the identity on the spectator and on the
 constants of the other coordinates,
 \begin{align*}
  \mathrm{Sym}(T'-T)&=T^\circ_j\otimes S^\circ_k+S^\circ_j\otimes T^\circ_k
   +S^\circ_j\otimes S^\circ_k+A^\circ_j\otimes A^\circ_k,\\
  \mathrm{Skew}(T'-T)&=(T_j+S_j)^\circ\otimes A^\circ_k
   +A^\circ_j\otimes(T_k+S_k)^\circ;
 \end{align*}
 in particular the skew part of $T'-T$ on $H_{\{j,k\}}$ has norm at most
 $\|(T_j+S_j)^\circ\|\,\|A^\circ_k\|+\|A^\circ_j\|\,\|(T_k+S_k)^\circ\|$.
\end{enumerate}
\end{lemma}

\begin{proof}
\emph{The setting.} For $S\ne S'$ pick $j$ in the symmetric difference; on pure
tensors the factors at $j$ are $\one$ and a function of mean zero, so
$H_S\perp H_{S'}$. Expanding $L^2(\pi_j)=\C\one\oplus L^2_0(\pi_j)$ for
$j\in\mathcal K$ shows that the sectors sum to $L^2(\pi)$ and that the functions
of $(\omega_0,\omega_W)$ form $\bigoplus_{S\subseteq W}H_S$.

(a) On a pure tensor $f_0\otimes\bigotimes_{j\in S}f_j\otimes\bigotimes_{j\notin S}\one$
with $f_j\in L^2_0(\pi_j)$, $T$ acts as $T_0f_0\otimes\bigotimes_{j\in S}T_j^\circ
f_j\otimes\bigotimes_{j\notin S}\one\in H_S(\pi')$, and such tensors span $H_S(\pi)$;
tensoring with $\C\one$ is isometric, and the norm of a tensor product of
operators is the product of the norms. The adjoint
$T^*=T_0^*\otimes\bigotimes_jT_j^*$ has the same form, each $T_j^*$ being
sector-preserving with $(T_j^*)^\circ=(T_j^\circ)^*$. A combination $R$ is the
orthogonal sum of its restrictions to the sectors, which are orthogonal on both
sides; so is an operator matrix with such entries, on the spaces
$H_S(\pi)\otimes\C^n$, and its quadratic form is the sum of theirs.

(b) By (a), $T'-T$ is $T_0\otimes(\bigotimes_{j\in S}T'^\circ_j-\bigotimes_{j\in
S}T^\circ_j)\otimes1$ on $H_S$, and $0$ on $H_\emptyset$, where both act as
$T_0\otimes1$. Telescoping over an ordering $j_1,\dots,j_s$ of $S$,
\[
 \bigotimes_iT'^\circ_{j_i}-\bigotimes_iT^\circ_{j_i}=\sum_{i=1}^s
 T'^\circ_{j_1}\otimes\dots\otimes T'^\circ_{j_{i-1}}\otimes
 (T'^\circ_{j_i}-T^\circ_{j_i})\otimes T^\circ_{j_{i+1}}\otimes\dots\otimes T^\circ_{j_s},
\]
of norm at most $\sum_id_{j_i}\prod_{k\ne j_i}\kappa_k$; the triangle inequality
gives the second bound. With $T_0=I$ and $\pi'=\pi$, $T'-T$ and its adjoint
preserve every $H_S$, so its symmetric and skew parts restrict to those of the
restrictions; on $H_{\{j\}}$ the restriction is $T'^\circ_j-T^\circ_j$. If every
$T_j$ is self-adjoint, so is $T|_{H_S}$, and the skew part of $(T'-T)|_{H_S}$ is
that of $\bigotimes_{j\in S}T'^\circ_j$, of norm at most $\delta'_S$. For
$s=|S|\ge2$, $sd\kappa^{s-1}=\kappa d\,s\kappa^{s-2}\le\kappa d\,s2^{2-s}\le
2d\kappa\le d$, and the three displayed bounds are the maxima over the sectors.
For the final assertion of (b), (a) applies to the combinations $T'-T$,
$\mathrm{Sym}(T'-T)$ and $\mathrm{Skew}(T'-T)$, whose restrictions to $H_S$ are
$(T'-T)|_{H_S}$ and its two parts. On $H_S$ with $|S|\ge2$ the second bound of
the display gives $\delta_S+\delta'_S$, and the symmetric and skew parts of an
operator have norm at most its own; for $j\ne k$ in $S$ every factor of
$\delta_S$ and of $\delta'_S$ other than those at $j$ and $k$ is at most one.

(c) $\Jproj=\bigotimes_j\Jproj_j$, with $\Jproj_jf=(\int f\,d\pi_j)\one$, is a
tensor product of sector-preserving operators with $\Jproj_j^\circ=0$; so it
agrees with $T$ on $H_\emptyset$ and vanishes on the other sectors, and
$T-\Jproj$ vanishes on $H_\emptyset$ and has norm
$\delta_S\le\min_{j\in S}\delta_j$ on $H_S$, $S\ne\emptyset$, with equality at
singletons; by (a) its norm is $\max_j\delta_j$. If each $T_j$ is self-adjoint
with spectrum in $[-c_j,1]$, then on $H_S$ the operator
$\bigotimes_{j\in S}T_j^\circ$ has eigenvalues $\prod_{j\in S}\lambda_j$,
$\lambda_j\in[-c_j,1]$; a negative product has a negative factor
$\lambda_i\ge-c_i$ and other factors of modulus at most one, so it is at least
$-\max_jc_j$. On $H_\emptyset$ the operator is the identity.

(d) For $X_S\in H_S\otimes E$, $X_{S'}\in H_{S'}\otimes E$ and $S\ne S'$,
$\int M(X_S,X_{S'})\,d\pi=\sum_{a,b}M_{ab}\int\overline{X_{S,a}}X_{S',b}\,d\pi=0$
in a basis of $E$. Hence
$\|\E[X\mid\omega_0,\omega_W]\|_M^2=\sum_{S\subseteq W}\|X_S\|_M^2$ for fixed $W$;
averaging over $\Jc$ gives $\sum_S\Pr(S\subseteq\Jc)\|X_S\|_M^2$, and
$\Pr(S\subseteq\Jc)=\delta_S$.

(e) The images $TX_S$ lie in the orthogonal sectors $H_S(\pi')$, and
$\|TX_S\|\le\|T_0\|\delta_S\|X_S\|$ by (a).

(f) With $\varphi_j=\one+c_j$, $c_j\in L^2_0(\pi_j)$,
$\bigotimes_j\varphi_j-\one=\sum_{S\ne\emptyset}\bigotimes_{j\in S}c_j\otimes
\bigotimes_{j\notin S}\one$, the term of $S$ lying in $H_S$; by orthogonality
the squared norm is $\sum_{S\ne\emptyset}\prod_{j\in S}\|c_j\|^2=
\prod_j(1+\|c_j\|^2)-1$.

(g) $\bigotimes_iA_i$ is self-adjoint with an orthonormal eigenbasis of product
vectors and eigenvalues $\prod_i\lambda_i$, $\lambda_i\in[-c_i,1+b_i]$, so
$|\lambda_i|\le1+b_i$. A product is at most $\prod_i(1+b_i)$, and a negative one
has a negative factor $\lambda_i\ge-c_i$, so it is at least
$-c_i\prod_{k\ne i}(1+b_k)$. Finally $1+b\le e^b$.

(h) As $T_j$ and $T'_j$ are sector-preserving, $T'_j-T_j$ and its adjoint
vanish on the constants, hence so do $S_j$, $A_j$ and their adjoints; and
$\ip\one{S_jf}=\ip{S_j^*\one}f=0$, likewise for $A_j$, so both map
$L^2_0(\pi_j)$ into itself. Thus $T'^\circ_j=T^\circ_j+S^\circ_j+A^\circ_j$ on
$L^2_0(\pi_j)$, with $T^\circ_j$, $S^\circ_j$ self-adjoint and $A^\circ_j$ skew. By
(a), $T'-T$ acts on $H_{\{j,k\}}$ as $T'^\circ_j\otimes T'^\circ_k-T^\circ_j\otimes
T^\circ_k$, with the identity on the other factors. Put
$U_j=(T_j+S_j)^\circ$ and $U_k=(T_k+S_k)^\circ$, both self-adjoint. The adjoint
of $(U_j+A^\circ_j)\otimes(U_k+A^\circ_k)$ is $(U_j-A^\circ_j)\otimes(U_k-A^\circ_k)$,
so its symmetric part is $U_j\otimes U_k+A^\circ_j\otimes A^\circ_k$ and its skew
part is $U_j\otimes A^\circ_k+A^\circ_j\otimes U_k$. The operator
$T^\circ_j\otimes T^\circ_k$ is self-adjoint; subtracting it and expanding
$U_j\otimes U_k-T^\circ_j\otimes T^\circ_k$ gives the two identities. The norm
bound follows from the triangle inequality and $\|X\otimes Y\|=\|X\|\,\|Y\|$.
\end{proof}

\begin{remark}[How the sector lemma is used]\label{rem:sector-uses}
Part (c) is why the comparison records of Part~V keep their actual marginals:
Markov operators of couplings with exact marginals, taken in those marginals,
have centred norm equal to the \emph{maximum} of the centred norms of their
factors under tensor products, not the sum or product of those norms. The maximal correlation of a product of couplings is the
largest factor correlation. For self-adjoint Markov operators the product
is bounded below by the minimum of zero and the smallest eigenvalues of
its factors; if at least one factor has a negative eigenvalue, the smallest
product eigenvalue equals the smallest factor eigenvalue. Operators normalized in a
reference space are not Markov; when their factor spectra lie in
$[-c_i,1+b_i]$, part (g) bounds every product eigenvalue above by
$\prod_i(1+b_i)$ and below by $-\max_i c_i\prod_{k\ne i}(1+b_k)$
(\cref{lem:spec:tensorcomp}). Part (b) prices a block difference: on a support
of at least two primes the difference of two products is of second order, and
when the old operator is self-adjoint only the new one contributes to the skew
part. Parts (a), (d) and (e) are the whole content of the index estimates
(\cref{cor:index-support-frame}(i), \cref{lem:record:common-mean}) and of
conductor decay; part
(f) says that for a product density $\bigotimes_j(\one+c_j)$ the
squared $L^2$ distance from $\one$ is
$\prod_j(1+\|c_j\|^2)-1$. The closing
clause of (b) prices a block whose prime differences have separate symmetric
and skew coefficients (the upper comparisons of \cref{sec:upper-comparison};
for the signed fold records of \cref{sec:middle-records}, which are
combinations of products, \cref{rem:signed-mixture-sectors} gives the same
decomposition), and (h) is the two-prime step of the matching records, whose
reference factor is self-adjoint (\cref{lem:oct24:matching-star}).
\end{remark}

\begin{remark}[Signed mixtures of products]\label{rem:signed-mixture-sectors}
Let the spectator be trivial, let $T_{i,j}$ ($i$ in a finite index set,
$j\in\mathcal K$) be sector-preserving, and put
$K=\sum_ic_i\bigotimes_{j\in\mathcal K}T_{i,j}$ with real $c_i$,
$\sum_ic_i=1$. By (a), $K$ maps every sector into the corresponding sector: it
is the identity on $H_\emptyset$, it acts on $H_{\{j\}}$ as the combination
$\sum_ic_iT^\circ_{i,j}$ at the single coordinate $j$, and on $H_S$ its norm is
at most $\sum_i|c_i|\prod_{j\in S}\|T^\circ_{i,j}\|$. The difference of two
such combinations decomposes in the same way: it vanishes on $H_\emptyset$, it
is the difference of the two coordinate combinations on $H_{\{j\}}$, and on
$H_S$ its norm is at most the sum of the two bounds; when $\pi'=\pi$, its
symmetric and skew parts restrict to the sectors in the same way, with the same
bound on $H_S$, and all these norms are maxima over the sectors. If the
$T_{i,j}$ are the operators of couplings $\pi_{i,j}$ of $\pi'_j$ and $\pi_j$,
then $K$ is the operator of the signed measure
$\sum_ic_i\bigotimes_j\pi_{i,j}$, which has the marginals $\pi'$ and $\pi$ and
total variation at most $\sum_i|c_i|$, whatever $|\mathcal K|$. A combination
taken at every coordinate, $\bigotimes_j(\sum_ic_iT_{i,j})$, has the same
marginals, but its total variation is the product of those of the coordinate
combinations, for which only the bound $(\sum_i|c_i|)^{|\mathcal K|}$ is
available. The fold records of the middle range are combinations of the first
kind, of three products with $\sum_i|c_i|=1+\tlaw$
(\cref{def:record:mixture}).
\end{remark}

\subsection{Fourier terms of large conductor are damped at the mixing primes \texorpdfstring{$p\ge\Pfive$}{p >= P\_5}}\label{subsec:lifts:decay}

A Fourier term supported on many prime coordinates is small after conditioning
on the other slots, provided each such prime supplies a mixing factor. The
orthogonality of distinct Fourier supports lets these factors multiply without
counting the supports. Reflection positivity is not an assumption in this
step.

\begin{definition}[Conditional mixing]\label{def:mixing}
Let $X$ be a finite set, $\rho$ a probability measure on $X^V$, and $Z\subseteq V$
a nonempty block of slots, with $\rho_Z$ the $\rho$-law of $z_Z$ and $\lambda$
the $\rho$-law of $z_{V\setminus Z}$. The \emph{conditional operator} from $Z$
to the other slots is
$T^{\rho,Z}f(y)=\E_\rho[f(z_Z)\mid z_{V\setminus Z}=y]$, from $L^2(\rho_Z)$ to
$L^2(\lambda)$; it fixes the constants, maps $\rho_Z$-mean-zero functions to
$\lambda$-mean-zero functions and is a contraction. For a prime $p$ and
$\gamma>0$, a probability measure $\rho$ on $\Fp^V$ is
\emph{$(\gamma,Z)$-mixing} if $\|T^{\rho,Z}f\|_{L^2(\lambda)}\le p^{-\gamma}
\|f\|_{L^2(\rho_Z)}$ for every $f$ with $\int f\,d\rho_Z=0$. For $Z=\{v\}$ we say
\emph{$(\gamma,v)$-mixing} (when $\rho$ has uniform one-slot marginals,
$\rho_{\{v\}}=\Unif(\Fp)$), and for a query pair $Z=\{u,v\}$, ordered as
$(u,v)$, \emph{$(\gamma,(u,v))$-mixing}.
\end{definition}

\begin{lemma}[Mixtures]\label{lem:mixture}
Let $Z$ and $Z'$ be finite sets and $\kappa=\sum_{i=1}^n\varpi_i\kappa_i$ with
$\varpi_i>0$, $\sum_i\varpi_i=1$, where the $\kappa_i$ are probability measures on
$Z\times Z'$ with the same $Z$-marginal $m$. Let $\kappa'$ and $\kappa'_i$ be the
$Z'$-marginals, and $T^\kappa f(z')=\E_\kappa[f(z)\mid z']$, $T^{\kappa_i}$
likewise, for $f\in L^2(m)$. Then $T^\kappa f=\sum_ih_iT^{\kappa_i}f$
$\kappa'$-almost everywhere, with $h_i=\varpi_i\,d\kappa'_i/d\kappa'$ and
$\sum_ih_i=1$, and
\[
 \|T^\kappa f\|^2_{L^2(\kappa')}\le\sum_i\varpi_i\|T^{\kappa_i}f\|^2_{L^2(\kappa'_i)} .
\]
Each $T^{\kappa_i}$, and $T^\kappa$, fixes the constants, is a contraction and
maps $m$-mean-zero functions to mean-zero functions. Consequently, if
$\|T^{\kappa_i}f\|\le c_i\|f\|$ for every $m$-mean-zero $f$, then
$\|T^\kappa f\|^2\le\sum_i\varpi_ic_i^2\|f\|^2$: a mixture of $(\gamma,Z)$-mixing
laws with a common $Z$-marginal is $(\gamma,Z)$-mixing, and a component of
weight $\varpi$ costs at most $\varpi^{1/2}$ in the norm.
\end{lemma}

\begin{proof}
For $z'$ with $\kappa'(z')>0$,
$T^\kappa f(z')=\sum_zf(z)\kappa(z,z')/\kappa'(z')=\sum_i(\varpi_i\kappa'_i(z')/
\kappa'(z'))\,T^{\kappa_i}f(z')$, the terms with $\kappa'_i(z')=0$ being absent.
By convexity $|T^\kappa f|^2\le\sum_ih_i|T^{\kappa_i}f|^2$ pointwise, and
$\int h_i|T^{\kappa_i}f|^2\,d\kappa'=\varpi_i\int|T^{\kappa_i}f|^2\,d\kappa'_i$.
Conditional Jensen gives $\|T^{\kappa_i}f\|^2_{L^2(\kappa'_i)}\le\int|f|^2\,dm$,
and $\int T^{\kappa_i}f\,d\kappa'_i=\int f\,dm$.
\end{proof}

The atoms below $\Pfive$, forming the finite low part of the functional,
have no quantitative mixing bound in (R4); in the lemma below they contribute
only the factor $C'^{\gamma}$. The conditional operator of the low part
has norm at most one, and, in the notation of the lemma, a frequency of
total conductor above $H$ has conductor above $H/C'$ at the tail primes
$p\ge\Pfive$.

\begin{lemma}[Conductor decay with a finite prefix]
\label{lem:record:prefix-conductor}
Let $X'$ be a finite product of cyclic groups of prime-power order, and let $C'$
be a positive integer divisible by the reduced denominator of every character
of $X'$. Let $\mathcal T$ be a finite set of primes not dividing $C'$. Let
$\rho'$ be a probability measure on $X'^V$ with uniform one-slot marginals. Fix a slot $v\in V$ and $\gamma>0$, and for $p\in\mathcal T$ let $\rho_p$ be a probability measure on $\Fp^V$ with uniform one-slot marginals that is $(\gamma,v)$-mixing at this same slot $v$ (\cref{def:mixing}). Put
$\rho=\rho'\otimes\bigotimes_{p\in\mathcal T}\rho_p$. Let
$g_w\colon X'\times X_{\mathcal T}\to\C$ $(w\in V)$ and $H\ge1$, and suppose that
$\widehat{g_v}(\xi)=0$ whenever $q(\xi)\le H$. Then
\[
 \Big|\int\prod_{w\in V}g_w(z_w)\,d\rho(z)\Big|
 \le\Big(\frac{C'}{H}\Big)^{\gamma}\|g_v\|_2\prod_{w\ne v}\|g_w\|_{2(d-1)},
\]
with norms for the uniform probability measure. If instead $\rho'=m\bar\rho'$
with $m\ge0$ and $\bar\rho'$ as above, the right-hand side is multiplied by $m$.
\end{lemma}

\begin{proof}
\emph{The conditional operator.} Let $\lambda'$, $\lambda_p$ be the laws of
$z'_{V\setminus\{v\}}$ under $\rho'$ and of $z_{p,V\setminus\{v\}}$ under $\rho_p$,
and $T'$, $T_p$ the conditional operators of $\rho'$, $\rho_p$ from slot $v$
(\cref{def:mixing}). By conditional Jensen and the uniform marginal of $z'_v$,
\begin{equation}\label{eq:pre:contraction}
 \|T'f\|_{L^2(\lambda')}^2=\E_{\rho'}\big|\E_{\rho'}[f(z'_v)\mid z'_{V\setminus\{v\}}]\big|^2
 \le\E_{\rho'}|f(z'_v)|^2=\|f\|_{L^2(X')}^2 ,
\end{equation}
so $\|T'\|\le1$; no mixing of $\rho'$ is used. Each $T_p$ fixes the constants and
$\int T_pf\,d\lambda_p=\E_{x\in\Fp}f(x)$, as $z_{p,v}$ is uniform; so $T_p$ is
sector-preserving, with $\|T_p^\circ\|\le p^{-\gamma}$ by mixing. The factors of
$\rho$ are independent, so the conditional law of $z_v$ given
$z_{V\setminus\{v\}}$ is the product of theirs, and the conditional operator of
$\rho$ from slot $v$ is $T=T'\otimes\bigotimes_{p\in\mathcal T}T_p$, from the
uniform space $X'\times X_{\mathcal T}$ to $L^2(\lambda)$,
$\lambda=\lambda'\otimes\bigotimes_p\lambda_p$.

\emph{The support condition.} Apply \cref{lem:sector} with the spectator $X'$,
$T_0=T'$, and the coordinates $p\in\mathcal T$. A character of
$X'\times X_{\mathcal T}$ is $\xi'+\xi_{\mathcal T}$, with $\xi'$ a character of
$X'$ and $\xi_{\mathcal T}$ of exact tail support $U\subseteq\mathcal T$; it lies in
$H_U$, and since $q(\xi')$ divides $C'$ and is coprime to
$q(\xi_{\mathcal T})=q_U$, its conductor is $q(\xi')q_U\le C'q_U$. The component
$g_{v,U}$ of $g_v$ in $H_U$ collects the Fourier terms of $g_v$ of tail support
$U$; if it is nonzero, then $H<q(\xi)\le C'q_U$ for some such term, so
$q_U>H/C'$. By \cref{lem:sector}(e),
\[
 \|Tg_v\|^2_{L^2(\lambda)}\le\sum_{U:\,q_U>H/C'}q_U^{-2\gamma}\|g_{v,U}\|_2^2
 \le\Big(\frac{C'}H\Big)^{2\gamma}\|g_v\|_2^2 .
\]

\emph{Conclusion.} Conditioning on $z_{V\setminus\{v\}}$, then applying
Cauchy--Schwarz in $L^2(\lambda)$ and H\"older with $d-1$ factors,
\[
 \Big|\int\prod_wg_w(z_w)\,d\rho\Big|=\Big|\int(Tg_v)(y)\prod_{w\ne v}g_w(y_w)\,d\lambda(y)\Big|
 \le\|Tg_v\|_{L^2(\lambda)}\prod_{w\ne v}\Big(\E_\lambda|g_w(y_w)|^{2(d-1)}\Big)^{\frac1{2(d-1)}} ,
\]
and the $\lambda$-law of $y_w$ is the $\rho$-law of $z_w$, uniform on
$X'\times X_{\mathcal T}$. The factor $m$ follows by multilinearity.
\end{proof}

With $X'$ trivial and $C'=1$ this is the same estimate with no low part.
The factor $C'^{\gamma}$ from the low part does not decay with $H$, and all
decay comes from the tail primes; the lemma is useful only because $C'$ is
bounded independently
of the scale. \Cref{sec:scale-comparison} applies it with $X'$ the product of the
groups of the atoms below $\Pfive$ that are present, $C'=Q_{\rm pre}$, and
$\mathcal T$ the tail primes $p\ge\Pfive$; \ref{it:record:R3} and
\ref{it:record:R4} supply its hypotheses, and no mixing is used below $\Pfive$.
The cancellation needs the same decay with a pair of slots as the source, and
a truncation of the pair supports.

\begin{lemma}[Decay by exact pair support]\label{lem:tr:pairsupport}
Let $\mathcal T$ be a finite set of primes, $(u,v)$ an ordered pair of distinct
slots and $\gamma>0$. For $p\in\mathcal T$ let $\rho_p$ be a probability measure
on $\Fp^V$ whose $(u,v)$-marginal is $\Unif(\Om)$ and which is
$(\gamma,(u,v))$-mixing, and put $\rho=\bigotimes_p\rho_p$. For $U\subseteq\mathcal T$
let $E_U$ be the conditional expectation, under $\bigotimes_p\Unif(\Om)$, onto
the pair coordinates at the primes of $U$, and $q_U=\prod_{p\in U}p$. Then for
$T_0\ge1$, every pair function $\varphi$ on $\prod_p\Om$ and every function $W$
of the labels of the other slots,
\[
 \Big|\int\varphi(z_u,z_v)W(z_{V\setminus\{u,v\}})\,d\rho(z)\Big|^2
 \le\|W\|_2^2\Big[T_0^{-2\gamma}\|\varphi\|_2^2
 +\sum_{q_U\le T_0}q_U^{-2\gamma}\|E_U\varphi\|_2^2\Big],
\]
with $\|\varphi\|_2$ and $\|E_U\varphi\|_2$ in $L^2(\bigotimes_p\Unif(\Om))$ and
$\|W\|_2$ in $L^2$ of the $\rho$-law $\lambda$ of the other slots.
\end{lemma}

\begin{proof}
The factors of $\rho$ are independent, so the conditional operator of $\rho$
from the pair $(u,v)$ to the other slots is $T=\bigotimes_pT_p$, with $T_p$
that of $\rho_p$;
each $T_p$ fixes the constants and integrates against $\lambda_p$ as against
$\Unif(\Om)$, so it is sector-preserving, with $\|T_p^\circ\|\le p^{-\gamma}$.
Conditioning on the other slots and Cauchy--Schwarz give
$|\int\varphi W\,d\rho|\le\|T\varphi\|_{L^2(\lambda)}\|W\|_2$, and
$\|T\varphi\|^2\le\sum_Uq_U^{-2\gamma}\|\varphi_U\|_2^2$ by \cref{lem:sector}(e),
with a trivial spectator. The
supports with $q_U>T_0$ contribute at most $T_0^{-2\gamma}\|\varphi\|^2$. For the
others, $E_U$ is the orthogonal projection onto $\bigoplus_{S\subseteq U}H_S$, so
$\varphi_U$ is the $U$-component of $E_U\varphi$ and
$\|\varphi_U\|_2\le\|E_U\varphi\|_2$.
\end{proof}

\paragraph{Consequences and what to read next.}
The slot array, its cuts and its six queries
(\cref{def:slot-array,lem:slot-array}), reflection positivity (\cref{def:rp}),
sorting sequences (\cref{def:sorting}) and the multilinear averages of
\cref{sec:positivity}, in the definitions only; the large sieve in Gram form and
Khintchine's inequality, both proved here. The reference laws, the comparison
law and admissibility
(\cref{hyp:record:reference,def:record:comparison-law,def:admissible}) are
interfaces whose instances are proved in Part~V.
